\section{Typed crossings}
\label{sec:rule}

\subsection{The argument}
\label{sec:argument}

Mathematicians apply a single symbol or predicate, such as $+$, $\le$, $\cong$,
\emph{compact} or \emph{valid}, to objects of quite different construction. On the
blackboard this abuse of notation is a feature. It compresses the exposition, because a
competent reader reconstructs which structure is meant. A proof assistant has no such
reader, and neither has the formal verification of a compiler or a kernel, nor the
evaluation of a model's output. There the following argument applies.

\begin{enumerate}[label=(P\arabic*)]
  \item Distinct structures may give one symbol different truth conditions. Subtraction
        on the natural numbers is not subtraction on the integers.
  \item Reusing the symbol is harmless where its truth conditions are invariant under the
        map between the structures. Where they are not, the reuse conceals a shift.
  \item Which of the two holds is a property of the map and the operation together. It
        cannot be read off the symbol.
  \item Where the properties of the map are not declared, every proof and every claim
        that crosses the boundary must establish them again, case by case. Where they are
        declared, the crossing can be automated and its failure conditions can be stated.
\end{enumerate}

\noindent
(P1)--(P3) are uncontroversial. (P4) is the practical premise, and the three cases are
offered as evidence for it.

\subsection{Vocabulary}
\label{sec:vocabulary}

The vocabulary is taken from Layered Event Modelling (LEM; \citealp{kreinzlem}), a method
proposal for layered models of cyber and operational risk. LEM's axioms are stipulations
adopted for their consequences, not theorems, and LEM says of itself that its transfer to
other domains is untested and would have to be argued, not read off a resemblance. The
cases below are such an argument for three domains. Only the following elements are
used.

\paragraph{Layers and locality.} A model assigns every node to exactly one layer (node
typing, LEM axiom A1). A predicate is defined at a declared layer and may be reused at
another only where its truth conditions and argument structure are demonstrably invariant
(predicate locality, A2). Where nodes of different layers refer to the same thing, that
co-reference is declared as a relation between distinct nodes (A3), so that a claim about
one is not silently transferred to the other.

\paragraph{Typed edges.} Every edge carries a declared type (A4). LEM types edges by what
would have to be true for the edge to fail, and distinguishes six types
(Table~\ref{tab:edges}). The table separates two things that the word ``fail'' can cover.
An edge does not hold when what it asserts is false under fixed conditions. An edge can
also depend on something declared, a scheme, a norm, a convention or a set of
assumptions, and can then cease to hold while the world stays as it is, because what
was declared is changed. That is a change of the conditions under which the edge is
judged, not an error in judging it. LEM uses this dependence as a test for
classificatory and normative edges. The table does not treat it as exclusive to them.
The type determines which properties are defined for the edge.
Duration, written $\Dt$, is defined for causal edges only: it is the duration of a
transition along which one event brings about the next. For every other type it is
undefined, not zero. The type of an edge is decided per edge, by test, and not by which
pair of layers the edge connects. Two of LEM's tests are used below: whether a non-zero
interval between the two ends is coherent, and whether an intervention can break the edge
without changing what its two ends are \citep{woodward2003}.

\begin{table}[t]
\centering\small
\begin{tabularx}{\textwidth}{@{}P{2.7cm}LLP{3.0cm}c@{}}
\toprule
Edge type & What the edge asserts & Does not hold when &
  Ceases to hold, the world unchanged, when & $\Dt$ \\
\midrule
Causal & One event brings about another, under intervention. &
  The transition does not occur, or is interrupted. & --- & defined \\
Constitutive & One state realises another at the same time; nothing propagates. &
  The realising states do not obtain. & --- & --- \\
Classificatory & An object satisfies the criteria of a declared scheme. &
  The criteria are not met. & The scheme is changed. & --- \\
Normative & An act or omission stands as fulfilment, or as violation, of a norm. &
  The asserted standing is wrong under the applicable norm. &
  The norm is amended, or no longer applies to the case. & --- \\
Representational & A representation corresponds to a referent. &
  The representation does not correspond (misrepresentation). &
  The convention under which it is read is changed. & --- \\
Evidential & An observation supports an inference. &
  The observation does not support the inference. &
  The background assumptions are changed. & --- \\
\bottomrule
\end{tabularx}
\caption{The six edge types of LEM, restated. The third column gives what makes the
asserted relation false under fixed conditions. The fourth gives what the edge depends
on that is declared and could be changed. The entries for the last two rows are this
paper's reading, not LEM's. Duration is defined for causal edges only.}
\label{tab:edges}
\end{table}

\paragraph{Five structures.} LEM separates five structures that layered models tend to
draw as one stack: a substrate (S1: the physical and computational entities that bear
everything else), an event abstraction (S2: events and the causal transitions between
them), representation (S3: descriptions and what they describe), epistemic access (S4:
observation and inference) and normative standing (S5: norms and the standing of acts
against them). They are not levels of one ladder. Sections~\ref{sec:refinement}
and~\ref{sec:reasoning} use them to sort the questions that one overloaded relation or
word leaves open.

\medskip
LEM has two further axioms, on the direction of causal edges across layers (A5) and on
the separation of type-level from token-level structure (A6). They are part of the
mechanisation (Section~\ref{sec:mech}) and play no role in the argument.

\subsection{The rule}
\label{sec:definition}

\begin{definition}[Crossing]
A \emph{crossing} is a map or relation $c$ between two structures $A$ and $B$ along which
a claim established in one is reused in the other.
\end{definition}

\begin{definition}[Typed crossing]
\label{def:typed}
A crossing is \emph{typed} when it declares
\begin{enumerate}[label=(\roman*)]
  \item its \emph{kind}: the role in which the crossing is used, which fixes what a
        failure of the crossing consists in; and
  \item what it \emph{carries}: for each predicate or operation concerned, whether and
        under which side conditions the crossing preserves it, reflects it, or both.
\end{enumerate}
\end{definition}

\noindent
For a crossing $c \subseteq A \times B$ and predicates $P_A$ on $A$ and $P_B$ on $B$,
$c$ \emph{preserves} $P$ when $c(a,b)$ and $P_A(a)$ imply $P_B(b)$, and \emph{reflects}
$P$ when $c(a,b)$ and $P_B(b)$ imply $P_A(a)$. Where both hold, $P$ is \emph{invariant}
under $c$. Which of these a reuse needs depends on the side on which the claim was
established. A property proved of a specification and asserted of an implementation
needs one direction only. The invariance asked for by predicate locality is the
two-directional case. Below, ``what a crossing preserves'' is short for this list. The
definition is formalised in Section~\ref{sec:mech-crossing}.

The rule proposed is that a crossing along which claims are carried be typed in this
sense. The two components answer different questions. The kind says what would count as
the crossing going wrong, and with it what sort of evidence can test it. The second
component says which claims may be carried without further proof.

The kind is a matter of role and not of mathematical structure. Whether a map is
injective, or a homomorphism for some operation, belongs to the second component and can
be proved. In which role the map is used, for instance whether its source is answerable
to its target, cannot be proved from the map and has to be declared;
Section~\ref{sec:three-kinds} gives the example.

``Structure'' is meant broadly: a type in a proof assistant, a state space with its
transition relation, a layer of a model. LEM's edge types are kinds in the sense of the
definition for layered event models. Section~\ref{sec:three-kinds} distinguishes three
roles for maps between types in a proof assistant. The two lists overlap, and the
overlap is offered as a parallel, not as an identity.

Table~\ref{tab:cases} summarises how the three cases instantiate the rule.

\begin{table}[t]
\centering\small
\begin{tabularx}{\textwidth}{@{}P{2.1cm}P{2.5cm}P{2.9cm}LL@{}}
\toprule
Case & Claim carried & Crossing & Declared locally & Left undeclared \\
\midrule
Casts in Lean~4 (Section~\ref{sec:casts}) &
  An identity or relation on $\N$ &
  The cast $\N \to \Z$ &
  Bundled morphisms; move, elim and squash lemmas &
  The kind of the coercion; a coercion instance says only that a canonical map exists \\
Refinement (Section~\ref{sec:refinement}) &
  A property of the specification &
  The simulation relation between implementation and specification &
  Preservation of properties of single executions; further properties by separate
  proofs &
  Which other properties survive; what ``verified'' is predicated of \\
Written reasoning (Section~\ref{sec:reasoning}) &
  ``The reasoning is correct'' &
  From the trace to the rules, and from the trace to the computation &
  Faithfulness distinguished from plausibility; tests by intervention &
  Which of the two crossings a given check tests \\
\bottomrule
\end{tabularx}
\caption{The three cases as instances of the rule.}
\label{tab:cases}
\end{table}
